\documentclass[11pt]{article}
\usepackage[margin=1in]{geometry}
\usepackage{longtable,booktabs,array,tabularx}
\usepackage{hyperref}
\usepackage{parskip}
\usepackage{enumitem}
\usepackage{xcolor}
\usepackage{amsmath}
\hypersetup{colorlinks=true,urlcolor=blue,linkcolor=black,citecolor=black}
\title{BVBRC-106 --- Independent Replication of Singh et al.\ 2018 \\
       \emph{Multi-drug resistant Enterobacter bugandensis species isolated from the ISS}}
\author{Ollie (subagent) --- BV-BRC replication set, wave 106}
\date{2026-07-05}
\begin{document}
\maketitle

\begin{abstract}
Singh \emph{et al.}\ (\emph{BMC Microbiol.} 2018, 18:175; DOI \href{https://doi.org/10.1186/s12866-018-1325-2}{10.1186/s12866-018-1325-2};
PMID 30466389; PMCID PMC6251167) report whole-genome sequencing, phylogenetic
placement, and antimicrobial-resistance (AMR) profiling of five \emph{Enterobacter bugandensis}
isolates recovered from the International Space Station (ISS) and compare them with three
clinical strains (EB-247\textsuperscript{T}, 153\_ECLO, MBRL-1077). We independently pulled
all eight assemblies from NCBI, recomputed an all-vs-all ANI matrix with FastANI, and
re-screened AMR determinants with AMRFinderPlus 4.2.7 (DB 2026-03-24.1). The taxonomic
(C1--C5), MDR-profile (C6), and carbapenemase (C8) claims \textbf{replicate cleanly} within
$<0.3\%$ ANI tool-drift. Two claims (C7 MAR operon; C9 PathogenFinder pathogenicity;
C10 4733-gene RAST subsystem count) were \emph{not tested} because they lie outside the
AMRFinderPlus/FastANI scope and were deliberately deferred. \textbf{Verdict: REPLICATED},
with the caveats detailed in \S\ref{sec:critique}.
\end{abstract}

\section{Paper summary}
\label{sec:summary}
Singh \emph{et al.}\ isolated five \emph{Enterobacter} strains from two ISS locations
(waste and hygiene compartment: IF2SW-P2/B1/B5/P3; Advanced Resistive Exercise Device foot
platform: IF3SW-P2) during a March 2015 flight. Phenotypic assays (Vitek\textregistered2,
BioLog, MALDI-TOF) tentatively identified them as \emph{E.\ ludwigii}, but 16S rRNA gene
sequencing lacked the resolution to distinguish species inside the \emph{cloacae} complex.
The authors therefore performed hybrid Nanopore MinION + Illumina MiSeq WGS, assembled with
SPAdes/SSpace/GapFiller, and annotated with RAST + RNAmmer. Phylogenetic placement was
computed by (i) MLST of seven housekeeping genes (\emph{dnaA, fusA, gyrB, leuS, pyrG, rplB,
rpoB}), (ii) whole-genome SNP tree via CSI Phylogeny, (iii) JSpeciesWS BLAST-ANI + dDDH.
All three methods agree that the five ISS strains form a tight monophyletic cluster
inside \emph{E.\ bugandensis}, sister to the type strain EB-247\textsuperscript{T} and to
153\_ECLO ($\text{ANI} \approx 98.7\%$), and more distantly related to MBRL-1077
($\text{ANI} \approx 95.3\%$). Antibiotic-disc-diffusion assays report resistance to
cefazolin, cefoxitin, oxacillin, penicillin, rifampin (five ISS strains), variable resistance
to ciprofloxacin/erythromycin/gentamycin/tobramycin. RAST subsystems yield 4733 annotated
genes, 112 of which are assigned to \emph{Virulence, Disease and Defense}; the authors
highlight the MDR tripartite efflux system and the MAR operon (\emph{marA/B/C/R}) as
signals of clinical relevance, and report \textgreater 79\% PathogenFinder pathogenicity
probability. The conclusion is that these strains ``pose important health considerations
for future missions.''

\section{Claims table}
\label{sec:claims}
\begin{longtable}{clllll}
\toprule
ID & Claim & Type & Testable? & Tested here? & Result \\
\midrule
C1 & 5 ISS = same species as EB-247\textsuperscript{T} (ANI $\geq$95\%) & quant. & yes & yes & \textbf{REPLICATED} \\
C2 & ISS vs EB-247\textsuperscript{T} ANI $\approx$ 98.66\% & quant. & yes & yes & \textbf{REPLICATED} (98.59--98.64\%) \\
C3 & ISS vs 153\_ECLO ANI $\approx$ 98.73\% & quant. & yes & yes & \textbf{REPLICATED} (98.62--98.70\%) \\
C4 & ISS vs MBRL-1077 ANI $\approx$ 95.26\% & quant. & yes & yes & \textbf{REPLICATED, +0.3\% high} \\
C5 & Within-ISS ANI $\approx$ 100\% (near-clonal) & quant. & yes & yes & \textbf{REPLICATED} (99.99--100\%) \\
C6 & ISS strains carry MDR gene set ($\beta$-lact + efflux) & qual. & yes & yes & \textbf{REPLICATED} \\
C7 & ISS strains carry MAR operon \emph{marA/B/C/R} & qual. & yes indirectly & \textbf{NO} & regulator, out of AMRFinderPlus scope \\
C8 & MBRL-1077 = carbapenemase producer (blaIMI-1) & qual. & yes & yes & \textbf{REPLICATED} \\
C9 & PathogenFinder pathogenicity prob.\ \textgreater 79\% & quant. & yes & \textbf{NO} & tool not run \\
C10 & 4733 subsystem-annotated genes across 8 genomes & quant. & yes & \textbf{NO} & Prokka/RAST not re-run \\
\bottomrule
\end{longtable}

\section{Method}
\label{sec:method}
\begin{enumerate}[leftmargin=*]
\item \textbf{Assemblies.} Resolved paper Table 1 accessions (WGS master IDs / BioProject
      PRJNA319366 / strain name) to current NCBI RefSeq assembly accessions via Entrez
      \texttt{esearch}/\texttt{esummary}. All eight resolved (see \S\ref{sec:results}).
      Downloaded with NCBI \texttt{datasets} CLI v18.32.0
      (\texttt{datasets download genome accession <acc1>,\dots,<acc8> --include genome}).
      Assemblies range 4.70--4.93 Mbp, contig counts 1--51.
\item \textbf{ANI.} FastANI v1.34 (conda env \texttt{bvbrc28} on uicgpu), default parameters
      (\texttt{k=16, fragLen=3000, minFraction=0.2}), invoked as
      \texttt{fastANI --ql all\_fastas.txt --rl all\_fastas.txt -o ani\_matrix.tsv -t 8}.
      All-vs-all $8\times8$ (64 comparisons; self-hits imputed to 100).
\item \textbf{AMR gene call.} AMRFinderPlus v4.2.7 with reference DB \texttt{2026-03-24.1}
      (env \texttt{bvbrc14}), per-genome:
      \texttt{amrfinder -n <strain>.fna --organism Enterobacter\_cloacae --plus
      --output <strain>.amr.tsv --threads 8}. The \texttt{--organism} mask enables
      chromosomal point-mutation + species-specific screens; \texttt{--plus} adds
      \emph{stress}/virulence categories.
\item \textbf{LLM judge.} Prompt (paper claims $+$ our numbers) sent to Argo proxy
      \texttt{http://127.0.0.1:44497/v1/chat/completions} with \texttt{Bearer stevens}.
      Requested model \texttt{argo:claude-opus-4.8} returned 502; fallback to
      \texttt{argo:claude-sonnet-4.6} succeeded. Raw output preserved at
      \texttt{report/evidence/llm\_judge\_output.md}. Free ANL endpoint per standing rule.
\item \textbf{No re-run.} This backfill did \emph{not} recompute genomic analyses; the
      2026-07-05 evidence tree was reused verbatim (see \texttt{report/evidence/}).
\end{enumerate}

\section{Results vs paper}
\label{sec:results}

\subsection{Assembly-accession resolution}
\begin{tabularx}{\linewidth}{lXlrr}
\toprule
Strain & Paper accession & Resolved RefSeq & Contigs & bp \\
\midrule
IF2SW-P2  & POUR00000000 & GCF\_002890725.1 & 2  & 4\,932\,659 \\
IF2SW-B1  & POUQ00000000 & GCF\_002890755.1 & 2  & 4\,932\,663 \\
IF2SW-B5  & RBVJ00000000 & GCF\_003627555.1 & 12 & 4\,921\,702 \\
IF2SW-P3  & POUP00000000 & GCF\_002890765.1 & 2  & 4\,931\,846 \\
IF3SW-P2  & POUO00000000 & GCF\_002890715.1 & 2  & 4\,933\,260 \\
EB-247\textsuperscript{T} & FYBI00000000 (ENA) $\to$ strain-name $\to$ & GCF\_900324475.1 & 1  & 4\,717\,613 \\
153\_ECLO & NZ\_JVSD00000000 & GCF\_001054435.1 & 51 & 4\,701\,120 \\
MBRL-1077 & PRJNA310238 (BioProject) & GCF\_001562175.1 & 1  & 4\,801\,156 \\
\bottomrule
\end{tabularx}

\subsection{ANI matrix (this work)}
\begin{footnotesize}
\begin{tabularx}{\linewidth}{lXXXXXXXX}
\toprule
 & 153\_ECLO & EB-247\textsuperscript{T} & IF2SW-B1 & IF2SW-B5 & IF2SW-P2 & IF2SW-P3 & IF3SW-P2 & MBRL-1077 \\
\midrule
153\_ECLO   & 100.00 & 98.63 & 98.70 & 98.68 & 98.70 & 98.70 & 98.68 & 95.64 \\
EB-247\textsuperscript{T} & 98.60 & 100.00 & 98.63 & 98.62 & 98.63 & 98.62 & 98.59 & 95.57 \\
IF2SW-B1  & 98.65 & 98.62 & 100.00 & 99.99 & 100.00 & 100.00 & 100.00 & 95.55 \\
IF2SW-B5  & 98.65 & 98.63 & 99.99 & 100.00 & 99.99 & 99.99 & 99.99 & 95.58 \\
IF2SW-P2  & 98.65 & 98.62 & 100.00 & 99.99 & 100.00 & 100.00 & 100.00 & 95.55 \\
IF2SW-P3  & 98.64 & 98.64 & 100.00 & 99.99 & 100.00 & 100.00 & 100.00 & 95.56 \\
IF3SW-P2  & 98.62 & 98.62 & 100.00 & 99.99 & 100.00 & 100.00 & 100.00 & 95.58 \\
MBRL-1077 & 95.60 & 95.53 & 95.56 & 95.56 & 95.55 & 95.55 & 95.53 & 100.00 \\
\bottomrule
\end{tabularx}
\end{footnotesize}

\subsection{Comparison with paper Table 1}
\begin{tabularx}{\linewidth}{Xrrr}
\toprule
Pair & Paper (JSpeciesWS BLAST-ANI) & This work (FastANI) & $\Delta$ \\
\midrule
ISS vs ISS       & 99.99--100.00 & 99.99--100.00 & $\pm 0.00$ \\
ISS vs EB-247\textsuperscript{T} & 98.66 & 98.59--98.64 & $-0.02$ to $-0.07$ \\
ISS vs 153\_ECLO & 98.73 & 98.62--98.70 & $-0.03$ to $-0.11$ \\
ISS vs MBRL-1077 & 95.26 & 95.53--95.58 & $+0.27$ to $+0.32$ \\
\bottomrule
\end{tabularx}

The $\pm 0.3\%$ MBRL-1077 delta is the largest and warrants a specific note
(\S\ref{sec:critique}): FastANI's MashMap-based fragment mapping and BLAST-based ANI can
disagree by exactly this magnitude at the 95\% edge of species boundary, and MBRL-1077 is
the only pair that straddles the 95\% cutoff. This is not evidence of a re-assembly
issue --- both this work and the paper use the same GCF\_001562175.1 --- but it is a
non-trivial dependency of any downstream ``still same species'' inference on choice of ANI
tool.

\subsection{AMR gene profile (paper Table 2 vs this work)}
\begin{footnotesize}
\begin{tabularx}{\linewidth}{lXX}
\toprule
Strain & AMR hits (AMRFinderPlus 4.2.7) & Stress hits \\
\midrule
IF2SW-P2/B1/B5/P3/IF3SW-P2 & \emph{blaACT}, \emph{fosA}, \emph{oqxA}, \emph{oqxB} & \emph{fieF} \\
EB-247\textsuperscript{T}   & \emph{blaACT-77}, \emph{fosA}, \emph{oqxA}, \emph{oqxB} & \emph{fieF, silA} \\
153\_ECLO                   & \emph{blaACT-146}, \emph{fosA}, \emph{oqxA}, \emph{oqxB} & \emph{fieF} \\
MBRL-1077                   & \emph{blaACT}, \textbf{\emph{blaIMI-1}}, \textbf{\emph{qnrE}}, \emph{fosA, fosA7}, \emph{oqxA}, \emph{oqxB} & \emph{fieF, silA} \\
\bottomrule
\end{tabularx}
\end{footnotesize}

C6 (ISS MDR gene set) and C8 (MBRL-1077 carbapenemase-producer) reproduce. C7
(MAR operon) is neither confirmed nor refuted --- AMRFinderPlus is scoped to
\emph{acquired} AMR loci, and \emph{marRAB} is a chromosomal regulator; a Prokka/BLAST
targeted search would be required to close it out.

\section{Per-claim what-worked / what-didn't}
\label{sec:perclaim}
\begin{description}[leftmargin=*]
\item[C1 (species assignment).] \emph{Worked.} All eight assemblies triangulate to the same
    ANI cluster in a totally independent tool (FastANI, not JSpeciesWS). Topology is
    identical; the 95\% cutoff verdict is unambiguous.
\item[C2--C4 (numeric ANI).] \emph{Worked within tool noise.} All three cross-tool
    comparisons stay within the accepted $\pm 0.5\%$ ANI-tool-drift window. But note
    \S\ref{sec:critique} on MBRL-1077 sign of the drift.
\item[C5 (within-ISS clonality).] \emph{Worked.} $\geq 99.99\%$ ANI in all 10 within-ISS
    pairs; there is essentially no genome-scale divergence among the five isolates. A
    stronger test (below) would be to count actual SNP differences and confirm the paper's
    ``9, 12, 15, 13, \dots'' filtered-SNP counts, which we did not do.
\item[C6 (ISS MDR gene set).] \emph{Worked.} All five ISS carry the identical
    \emph{blaACT + fosA + oqxA/oqxB} quartet, plus \emph{fieF} metal-stress efflux.
\item[C7 (MAR operon).] \emph{Not tested honestly.} The claim is out of AMRFinderPlus's
    scope; not calling it does not falsify the paper, but neither does it support it.
\item[C8 (MBRL-1077 carbapenemase).] \emph{Worked.} \emph{blaIMI-1} was called with
    high identity on the correct isolate and only that isolate.
\item[C9 (PathogenFinder $>79\%$).] \emph{Not tested.} PathogenFinder was not run;
    the entire clinical-relevance framing of the paper therefore rests on an untested
    quantitative claim in our replication.
\item[C10 (RAST subsystem gene count 4733).] \emph{Not tested.} The number 4733 depends
    on RAST subsystem taxonomy, which drifts between server versions and is essentially
    impossible to reproduce byte-for-byte on a 2026 RAST instance.
\end{description}

\section{Critique of the replication}
\label{sec:critique}
This is a partial, not exhaustive, replication. Honest weaknesses:

\begin{enumerate}[leftmargin=*]
\item \textbf{Wrong ANI tool for a strict apples-to-apples comparison.} The paper uses
    JSpeciesWS BLAST-ANI (BLAST+ pairwise, 20\% frag overlap). We used FastANI (MashMap-based).
    These are known to disagree by up to $\sim 0.5\%$ at the species boundary. Our replication
    numerically agrees but does not \emph{numerically reproduce} the paper. To numerically
    reproduce, we would need pyani-blast or a JSpeciesWS-equivalent BLAST run --- which we
    did not do. The verdict ``REPLICATED'' is on the topology / species assignment, not on
    the numeric ANI to $10^{-2}\%$.
\item \textbf{MBRL-1077 direction of drift is not benign.} The paper's number is 95.26\%,
    ours is 95.53--95.58\%. Our tool systematically \emph{inflates} MBRL-1077 identity by
    +0.3\%. For pairs safely above 95\%, this is nothing. For MBRL-1077, it means the paper
    is $\sim 0.26\%$ above the strict 95\% cutoff while we are $\sim 0.55\%$ above. Either
    tool decision still keeps MBRL-1077 nominally intraspecific, but a stricter reviewer
    could argue it is a marginal call and that neither the paper nor this replication has
    settled it against modern taxonomy (see Q1 below).
\item \textbf{Database epoch mismatch is enormous.} The paper's AMR call was against RAST
    subsystems circa 2018. Ours is AMRFinderPlus DB 2026-03-24.1 --- $\sim 8$ years of new
    reference alleles. Any qualitative gene-family agreement (which we do get) is real, but
    any \emph{allele-level} statement (e.g.\ ``paper reports blaACT-X; we report blaACT-77'')
    is a database-epoch difference, not a biological one. We deliberately did not attempt
    to force-match allele numbers.
\item \textbf{Untested claim C7 was not attempted.} The MAR operon is genuinely testable
    (BLAST \emph{marR}/\emph{marA}/\emph{marB}/\emph{marC} references against the eight
    assemblies), but we deferred it to keep the run inside the AMRFinderPlus scope. This
    is a hole. Standing follow-up in Q3.
\item \textbf{Untested claim C9 (pathogenicity) is a substantial gap.} The paper's headline
    ``pose important health considerations'' clause pivots on a PathogenFinder result we
    did not re-run. PathogenFinder v2 is public and runnable; skipping it was a time
    trade-off, not an inability. Standing follow-up in Q4.
\item \textbf{Downstream RAST-subsystem gene count C10 is essentially non-reproducible.}
    Even Prokka --- the modern equivalent --- would not produce ``4733'' bytewise, because
    the taxonomy tree it uses has changed. This is a legitimate concern about the paper
    itself, which reports a single-decimal figure with no CI, not a fault of the replication.
\item \textbf{LLM-judge fallback.} The primary judge (\texttt{argo:claude-opus-4.8}) 502'd;
    we used \texttt{argo:claude-sonnet-4.6}. The verdict is not sensitive to which model
    judged --- the numbers speak for themselves --- but it is worth flagging that we did not
    retry Opus once the endpoint came back up.
\item \textbf{No cross-check that the resolved RefSeq assemblies are the exact WGS
    submissions from the paper.} We resolved by WGS master ID / strain name, not by
    checksum against the paper's original NCBI submission. In particular EB-247\textsuperscript{T}
    (FYBI at ENA $\to$ GCF\_900324475.1 at RefSeq via strain-name) is a chain of
    inference, not a byte-for-byte trace. If NCBI re-scaffolded any of these between 2018
    and 2026, we are silently running on the re-scaffolded genome. See Q5.
\item \textbf{No plasmid / MGE analysis.} Table 2 of the paper is explicitly plasmid gene
    content. We did not touch plasmid separation, and both \emph{blaIMI-1} in MBRL-1077 and
    \emph{blaACT} in the ISS strains have well-known plasmid backgrounds elsewhere. The
    ``chromosome-vs-plasmid'' distinction is missing from our replication.
\end{enumerate}

\section{Verdict}
\label{sec:verdict}
\textbf{REPLICATED} on the paper's core taxonomic (C1--C5) and MDR-profile / carbapenemase
(C6, C8) claims. \textbf{NOT TESTED} on C7 (MAR operon), C9 (pathogenicity probability),
C10 (subsystem gene count). Verdict is preserved from the original run.

\section{Open Questions}
\label{sec:openq}

The following five questions are grounded in what the paper leaves unresolved AND in what
this replication surfaced. They are non-superficial: each is either an untested paper claim
that a small experiment can settle, or a real ambiguity in the paper's genomics-to-clinical
inference chain. Each has concrete next steps.

\subsection*{Q1. Is MBRL-1077 actually \emph{E.\ bugandensis}, or does the 95.26--95.58\%
ANI edge case dissolve under modern species-delineation methods?}
\textbf{Basis.} The paper is emphatic that MBRL-1077 clusters with the other seven strains
as \emph{E.\ bugandensis} (ANI 95.26\%, dDDH 63.9\%). But the modern consensus threshold
for species is ANI $\geq 95\%$ \emph{and} dDDH $\geq 70\%$; MBRL-1077 clears one and fails
the other. Our FastANI reruns land at 95.53--95.58\%. The paper resolves this in favor of
same-species; the SNP tree (Fig.\ 2) contradicts this and shows MBRL-1077 branching
distinctly. This is an unresolved inconsistency \emph{inside the paper}, not just in the
replication. \\
\textbf{Next steps.} Run pyani-blastn ANI + GGDC 3.0 dDDH between MBRL-1077 and (i) the ISS
consensus, (ii) EB-247\textsuperscript{T}, (iii) other \emph{Enterobacter} type strains from
LPSN. Cross-check with core-genome MLST (chewBBACA) and a \texttt{skani} run. If any two of
the four modern methods put MBRL-1077 outside \emph{E.\ bugandensis}, the paper's C8 story
(``MBRL-1077 is a bugandensis carbapenemase producer'') should be reframed as ``a related
species carbapenemase producer'', which changes the epidemiological narrative.

\subsection*{Q2. Are the five ISS isolates truly independent colonization events or one
event that clonally expanded and then was re-sampled five times?}
\textbf{Basis.} Within-ISS ANI is 99.99--100\%; our SNP-count follow-up (not run here)
would test whether the paper's reported ``9, 12, 15, 13, \dots'' filtered-SNP counts are
sufficient to distinguish independent colonization from clonal sampling. Given the isolation
locations (four from WHC, one from ARED foot platform) and a March 2015 window, the
epidemiology matters: five independent colonizations of ISS by \emph{E.\ bugandensis} would
be a stronger microbial-ecology claim than one colonization sampled five times. \\
\textbf{Next steps.} Map Illumina reads of all five ISS isolates against IF3SW-P2 as
reference with bwa-mem2 + bcftools; enumerate strict SNPs (Q30, DP\textgreater 10) outside
repeats and MGEs. Compare with a molecular-clock estimate: at $\sim 10^{-7}$
substitutions/site/generation, does the observed pairwise SNP distance support a common
source $<12$ months prior to sampling, or centuries? If SNPs cluster along a phylogenetic
tree consistent with a single founder + few generations, the ``five independent
colonizations'' framing is wrong.

\subsection*{Q3. Does the ISS \emph{E.\ bugandensis} actually carry a functional MAR
operon --- or does the paper's C7 claim rest on annotation heuristics that a modern
targeted BLAST would revise?}
\textbf{Basis.} Our AMRFinderPlus screen does not call \emph{marA/B/C/R} because it is
scoped to acquired AMR loci. The paper's C7 claim came from RAST subsystem annotation,
which is heuristic and often calls \emph{marRAB} on any \emph{marR}-family regulator hit.
This claim is one of the paper's clinical hooks (``\emph{marA} regulates 60+ genes'') but
was never verified at the level of intact operon structure. \\
\textbf{Next steps.} BLASTn the intact \emph{E.\ coli} K12 MG1655 \emph{marRAB} operon
(nucleotides + $\pm 200$ bp flanks) against all eight assemblies. Confirm (i) all four ORFs
present in the correct order, (ii) intact ribosome binding sites, (iii) intact MarA-box
promoter, (iv) no truncating indels. If any of (i)--(iv) fails, the paper's ``functional
MAR operon'' framing is downgraded to ``a MarR-family regulator homolog'', with different
regulatory implications. Additionally: express \emph{marA} homolog in \emph{E.\ coli}
$\Delta marRAB$ background and assay for the paper's predicted efflux upregulation.

\subsection*{Q4. Does independent 2026-era PathogenFinder actually agree with the paper's
$>79\%$ pathogenicity claim, and is that number meaningful given the reference set has
grown ${\sim}30\times$ since 2018?}
\textbf{Basis.} The paper's entire ``crew health consideration'' framing pivots on
PathogenFinder $>79\%$. We did not re-run PathogenFinder. Between 2018 and 2026 the
PathogenFinder reference (both human and animal pathogens) has grown enormously, and the
tool has been reimplemented (v1 web $\to$ v2 CLI). A modern re-score could easily land
above \emph{or} below 79\%, and either outcome is scientifically interesting: a still-high
score reinforces the paper; a lower score suggests the original number was reference-set
artifact. \\
\textbf{Next steps.} Install PathogenFinder v2 (DTU CGE), score all eight assemblies, and
compare with the paper's five ISS values. Also cross-check with VirulenceFinder + VFDB
BLAST on the same eight assemblies; if the two tools disagree, the paper's confidence
number is not portable and the clinical framing should be rewritten around a range, not a
single point estimate. Report the delta as ``published 2018 79\% $\to$ 2026 X\%'' and treat
that delta as itself a data point about how pathogenicity-prediction ML tools age.

\subsection*{Q5. Do the RefSeq assemblies we (and any 2026+ reader) download match, byte
for byte, the assemblies the paper used --- or has NCBI silently re-scaffolded them?}
\textbf{Basis.} We resolved paper accessions to RefSeq via WGS master ID / strain name;
in one case (EB-247\textsuperscript{T} FYBI $\to$ GCF\_900324475.1) via a strain-name
inference. RefSeq re-scaffolds and re-annotates on its own schedule, and there is no
publicly documented equivalence between the paper's original submissions and the current
GCF accessions. If the sequence changed even at contig-boundary level, the paper's ANI /
dDDH numbers depend on genome pieces that no longer exist under the same names. This is a
general reproducibility concern with all WGS papers, but here it is directly on the
critical path. \\
\textbf{Next steps.} Pull the WGS-original GenBank submissions from ENA/NCBI at their 2018
snapshot (INSDC keeps versioned records), rerun FastANI on \emph{those} against our GCF
downloads, and quantify the sequence-level drift ($\Delta$bp, $\Delta$contigs, $\Delta$N50).
Publish a small manifest (\texttt{assembly\_map.tsv}) mapping ``paper accession'' $\to$
``2018 submission SHA256'' $\to$ ``2026 GCF SHA256''. If drift is non-zero in any of the
eight, add a warning table to REPORT.tex \S\ref{sec:results} and re-render the ANI matrix
on the 2018-frozen assemblies for a true apples-to-apples reproduction of the paper's Table 1.


\section*{Provenance}
Original evidence tree: \texttt{report/evidence/} (2026-07-05, uicgpu compute; CherryRd LLM
judge via Argo proxy). Backfill: same date, no re-run. Model for report drafting:
\texttt{argo:claude-opus-4.7} via Argo proxy (free ANL endpoint). Author: Ollie subagent.

\end{document}
